% =============================================================================
% CompStat Rio — Guia de Leitura da Aplicação
% Autor: Arthur Lins  |  Compila com pdflatex (Overleaf-friendly).
% Figuras esperadas em ./figuras/ (geradas em prompt separado).
% =============================================================================
\documentclass[11pt,a4paper]{article}

\usepackage[utf8]{inputenc}
\usepackage[T1]{fontenc}
\usepackage[brazilian]{babel}
\usepackage[margin=2.5cm]{geometry}
\usepackage{graphicx}
\usepackage{amsmath}
\usepackage{amssymb}
\usepackage{booktabs}
\usepackage{caption}
\usepackage{subcaption}
\usepackage{longtable}
\usepackage{microtype}
\usepackage{xcolor}
\usepackage{float}  % habilita o especificador [H] para travar figuras na posição
\usepackage[
  colorlinks=true,
  linkcolor=blue!50!black,
  urlcolor=blue!60!black,
  citecolor=blue!60!black
]{hyperref}

\graphicspath{{figuras/}}

% --- Inclusão de figura com fallback (compila antes dos PNGs existirem) ---
% Usamos \detokenize no fallback para imprimir nomes com '_' literalmente,
% sem que o TeX interprete underscore como subscrito fora de modo matemático.
\newcommand{\figfile}[2]{%
  \IfFileExists{figuras/#1}{%
    \includegraphics[width=#2]{#1}%
  }{%
    \fbox{\parbox{#2}{\centering\small\itshape%
      [figura pendente: \texttt{\detokenize{#1}}]}}%
  }%
}

% --- Metadados ---
\title{%
  \textbf{CompStat Rio} \\[0.4em]
  \large{Guia de Leitura da Aplicação} \\[0.3em]
  \normalsize\itshape Inteligência territorial para segurança pública municipal
}
\author{Arthur Lins}
\date{\today}

\begin{document}

% --- Capa ---
\begin{titlepage}
  \centering
  \vspace*{2cm}
  {\Huge \textbf{CompStat Rio}} \\[0.6cm]
  {\Large Guia de Leitura da Aplicação} \\[0.4cm]
  {\large\itshape Inteligência territorial para segurança pública municipal} \\[2cm]

  {\large \textbf{Arthur Lins}} \\[0.3cm]
  {\large \today}
  \vfill
  \noindent\rule{\linewidth}{0.4pt}
  \begin{flushleft}
    \small
    Documento didático destinado à alta gestão pública. Apresenta o que a
    aplicação CompStat Rio entrega, como navegar nela, e a especificação
    econométrica por trás do mapa preditivo de risco. \\[0.3em]
    \emph{Decisão final humana. A plataforma produz rascunho auditável; a
    equipe valida antes de qualquer ação.}
  \end{flushleft}
  \noindent\rule{\linewidth}{0.4pt}
\end{titlepage}

\tableofcontents
\newpage

% =============================================================================
\section{Introdução}
\label{sec:introducao}

O \textbf{CompStat Rio} é a plataforma de inteligência territorial da
Prefeitura do Rio de Janeiro voltada para segurança pública municipal.
O modelo de gestão CompStat realiza reuniões periódicas em que o Prefeito
e a alta gestão decidem, com base em dados, onde a Força Municipal (FM)
deve atuar e quais fatores urbanos devem ser resolvidos pelos órgãos
responsáveis (Comlurb, RioLuz, Seconserva, CET-Rio, SEOP, SMAS, SMTR,
GM-Rio).

\paragraph{O problema dos silos.} As fontes de dados — ocorrências
georreferenciadas, denúncias do Disque Denúncia, relatórios de
inteligência (RELINTs), fatores urbanos mapeados em campo e câmeras —
vivem em sistemas isolados. Cruzar manualmente esses silos para produzir
um Relatório Analítico de Área consome horas que poderiam ser dedicadas à
análise. O CompStat Rio automatiza esse cruzamento.

\paragraph{Três camadas, uma plataforma.} A inteligência central
sobrepõe, no mesmo ponto do mapa, três camadas:
\begin{itemize}
  \item \textbf{Mancha criminal} — onde o crime acontece, na forma de
        ocorrências de roubo e furto georreferenciadas.
  \item \textbf{Fator urbano} — qual condição ambiental facilita o crime
        (iluminação deficiente, calçada obstruída, vegetação cobrindo
        postes, mobiliário abandonado, etc.). É a alavanca da
        intervenção municipal.
  \item \textbf{Dinâmica criminal} — como o crime acontece (modus
        operandi, rotas de fuga, perfis de vítima), inferida do texto
        livre de RELINTs e Disque Denúncia.
\end{itemize}

\paragraph{O papel da IA.} A inteligência artificial amplifica a análise
— sintetiza texto, cruza dados e propõe rascunhos — mas \emph{nunca
substitui o gestor}. Os guardrails da plataforma são detalhados na
seção~\ref{sec:guardrails}.

\begin{figure}[H]
  \centering
  \figfile{01_home_lista_areas.png}{0.95\linewidth}
  \caption{Página inicial do CompStat Rio: lista de áreas da Força
  Municipal ordenadas por urgência (volume de ocorrências de roubo e
  furto no período do piloto). Selecionar uma área abre o relatório
  analítico completo.}
  \label{fig:home}
\end{figure}

% =============================================================================
\section{Como navegar}
\label{sec:navegar}

A aplicação tem duas visões principais, alternadas por um seletor de
abas centralizado no topo:

\begin{itemize}
  \item \textbf{Panorama Operacional} — a entrada padrão: lista de áreas
        prioritárias e, dentro de cada uma, o relatório analítico
        completo.
  \item \textbf{Mapa Preditivo} — a previsão econométrica de risco para
        as próximas 1, 2 e 4 semanas (seção~\ref{sec:preditivo}). \textbf{\emph{Essa é a maior inovação do nosso produto, e a que temos mais esperança de que gere impacto real na produtividade de gestores públicos.}}
\end{itemize}

\begin{figure}[H]
  \centering
  \figfile{02_topbar_abas.png}{0.95\linewidth}
  \caption{Seletor de abas no topo. As duas visões têm igual destaque
  visual; a aba ativa fica em azul-petróleo.}
  \label{fig:topbar}
\end{figure}

\noindent A lista de áreas (figura~\ref{fig:home}) é ordenada por
ranking de volume de ocorrências. Ao clicar em uma área, o usuário entra
no \textbf{Relatório Analítico de Área}, que segue o formato dos
Relatórios de Inteligência (RELINTs) já adotados pela FM. Cada relatório
é estruturado em 10 seções (S1 a S10), descritas a seguir.

% =============================================================================
\section{Panorama Operacional: o Relatório Analítico de Área}
\label{sec:panorama}

\subsection{S1 — Identificação}
Cabeçalho da área: nome, polígono oficial da FM, bairros principais,
domínios territoriais de organizações criminosas identificados, AISP e
contagem de trechos críticos detectados pelo motor de cruzamento.

\begin{figure}[H]
  \centering
  \figfile{04_secao_s1_identificacao.png}{0.95\linewidth}
  \caption{Seção 1 — Identificação da área.}
  \label{fig:s1}
\end{figure}

\subsection{S2 — Mapa de Hotspots}
Visualização espacial integrada: heatmap de ocorrências sobreposto ao
polígono da área FM, com camadas de câmeras, fatores urbanos e segmentos
críticos coloridos por escore. Cada elemento tem legenda e popup.

\begin{figure}[H]
  \centering
  \figfile{05_secao_s2_mapa.png}{0.95\linewidth}
  \caption{Seção 2 — Mapa de hotspots.}
  \label{fig:s2}
\end{figure}

\subsection{S3 — Resumo Executivo}
Respostas geradas pela IA para as quatro perguntas operacionais do
CompStat: \emph{onde} patrulhar, \emph{quando} patrulhar, \emph{como}
empregar a FM e \emph{como} os órgãos devem resolver os fatores urbanos.
Cada resposta vem com \emph{provenance} (origem dos dados e nível de
confiança).

\begin{figure}[H]
  \centering
  \figfile{06_secao_s3_resumo.png}{0.95\linewidth}
  \caption{Seção 3 — Resumo executivo com proveniência.}
  \label{fig:s3}
\end{figure}

\subsection{S4 — Ocorrências no Período}
\label{ssec:s4}
Indicadores quantitativos de ocorrências de roubo registradas na área.
O período de cobertura está explícito no topo da seção: \textbf{janela
do piloto de jan/2023 a dez/2024}. Como não há período anterior
comparável nos dados disponíveis, a métrica de \emph{variação vs.\
período anterior} não é exibida (omissão honesta em vez de cálculo
inventado).

\begin{figure}[H]
  \centering
  \figfile{07_secao_s4_ocorrencias.png}{0.95\linewidth}
  \caption{Seção 4 — Indicadores de ocorrências e período do piloto.}
  \label{fig:s4}
\end{figure}

\subsection{S5 — Análise Temporal}
Mapa de calor 7$\times$24 (dia da semana $\times$ hora) que mostra os
picos de incidência na área. Acompanha uma síntese textual gerada pela
IA, com proveniência.

\begin{figure}[H]
  \centering
  \figfile{08_secao_s5_temporal.png}{0.95\linewidth}
  \caption{Seção 5 — Mapa de calor temporal.}
  \label{fig:s5}
\end{figure}

\subsection{S6 — Dinâmica Criminal}
Síntese qualitativa do modus operandi, rotas de fuga e padrões
observados, extraída do texto livre de RELINTs e Disque Denúncia. Cada
afirmação cita a fonte literal e o nível de confiança --- \emph{texto
livre é tratado como indício, não como fato}.

\begin{figure}[H]
  \centering
  \figfile{09_secao_s6_dinamica.png}{0.95\linewidth}
  \caption{Seção 6 — Dinâmica criminal com proveniência.}
  \label{fig:s6}
\end{figure}

\subsection{S7 — Emprego do Efetivo da FM}
\label{ssec:s7}
Sugestão de cobertura, horário e modalidade (a pé, viatura, moto) para a
Força Municipal. A sugestão é uma \textbf{recomendação da IA construída
a partir de dois insumos}: (i) o \emph{pico temporal} (dia e hora de
maior incidência) e (ii) o \emph{deslocamento do autor} --- um indício
extraído de RELINT e Disque Denúncia, tratado como inteligência, não
como fato. A sugestão \emph{não usa} o escore de prioridade da S10. A
decisão final é sempre do gestor.

\begin{figure}[H]
  \centering
  \figfile{10_secao_s7_efetivo.png}{0.95\linewidth}
  \caption{Seção 7 — Sugestão de emprego do efetivo, com explicação dos
  insumos usados pela IA.}
  \label{fig:s7}
\end{figure}

\subsection{S8 — Fatores Urbanos}
Inventário dos fatores ambientais mapeados na área (iluminação,
vegetação, calçada, lixo, mobiliário, situação de rua, trânsito),
agregados por tipo e órgão responsável \cite{ObrienSampson2015,
Wheeler2018}.

\begin{figure}[H]
  \centering
  \figfile{11_secao_s8_fatores.png}{0.95\linewidth}
  \caption{Seção 8 — Fatores urbanos por tipo e órgão.}
  \label{fig:s8}
\end{figure}

\subsection{S9 — Câmeras}
Cobertura de câmeras CIVITAS/COR na área, com identificação visual de
pontos cegos.

\begin{figure}[H]
  \centering
  \figfile{12_secao_s9_cameras.png}{0.95\linewidth}
  \caption{Seção 9 — Cobertura de câmeras.}
  \label{fig:s9}
\end{figure}

\subsection{S10 — Coincidências de Alto Risco e Plano de Ação}
\label{ssec:s10}
Esta é a seção onde o cruzamento das três camadas se materializa em
prioridade. Para cada segmento crítico, é calculado um \emph{escore}
$\in [0, 10]$:

\begin{equation}
\label{eq:score}
\text{score} = \min\!\Bigl(10\,;\;
  5\cdot \text{densidade}_{\text{norm}}
  + 2\cdot \mathbb{1}[\text{tem fator}]
  + 1{,}5\cdot \mathbb{1}[\text{tem dinâmica}]
  + 1{,}5\cdot \mathbb{1}[\text{lacuna câmera}]
\Bigr)
\end{equation}

\noindent\textbf{Lendo cada termo em linguagem simples:}
\begin{itemize}
  \item $\text{densidade}_{\text{norm}}$ é a densidade de ocorrências do
        trecho dividida pela densidade máxima da área (varia de 0 a 1).
        Peso~5: é o sinal mais forte porque é a evidência mais direta de
        que algo está acontecendo.
  \item $\mathbb{1}[\text{tem fator}]$ vale 1 se há fator urbano mapeado
        no raio do trecho, 0 caso contrário. Peso~2: a presença de uma
        causa removível eleva a prioridade.
  \item $\mathbb{1}[\text{tem dinâmica}]$ vale 1 se há indício
        qualitativo (RELINT ou Disque) sobre a área. Peso~1{,}5: texto
        livre é indício, peso menor reflete a incerteza.
  \item $\mathbb{1}[\text{lacuna câmera}]$ vale 1 se não há câmera no
        raio. Peso~1{,}5: lacuna de cobertura agrava o segmento.
\end{itemize}

\noindent Os \textbf{pesos são uma escolha do produto} (não foram
estimados por nenhum modelo) e estão extraídos em constantes nomeadas no
código (\texttt{app/backend/match/score.py}). A fórmula é exibida na
própria tela, no painel da seção S10, para que o gestor não trabalhe com
caixa-preta.

\begin{figure}[H]
  \centering
  \figfile{13_secao_s10_coincidencias.png}{0.95\linewidth}
  \caption{Seção 10 — Coincidências de alto risco. A nota acima do
  painel mostra a fórmula do escore com os pesos atuais.}
  \label{fig:s10}
\end{figure}

% =============================================================================
\section{Mapa Preditivo}
\label{sec:preditivo}

\subsection{O que é e como ler}
\label{ssec:pred_overview}

O Mapa Preditivo é a frente \emph{econométrica} da plataforma. Em vez de
um escore aditivo desenhado à mão, aqui um \textbf{modelo de regressão
logística} estima a probabilidade $\Pr(y_{i,t+s}=1)$ de \emph{haver pelo
menos um roubo} em um micropolígono $i$ na semana $t{+}s$, com
$s \in \{1, 2, 4\}$ semanas à frente.

A unidade espacial é o \textbf{hexágono H3 de resolução 9}
(aproximadamente 0{,}1 km$^2$) dentro do polígono da FM; a unidade
temporal é a \emph{semana ISO}. O método dialoga com a literatura de
\emph{near-repeat victimization} \cite{JohnsonBowers2004,
MohlerEtAl2011} e \emph{risk terrain modeling} \cite{CaplanKennedyMiller2011}.

\begin{figure}[H]
  \centering
  \figfile{14_mapa_preditivo_geral.png}{0.95\linewidth}
  \caption{Mapa Preditivo --- visão geral, com seletor de área e painel
  analítico em três abas (Drivers, Métricas, Coeficientes).}
  \label{fig:pred_overview}
\end{figure}

\subsection{Especificação econométrica}
\label{ssec:pred_spec}

O modelo é um logit em painel hex-semana, com efeitos fixos de área FM:

\begin{equation}
\label{eq:logit}
\Pr\!\left(y_{i,t+s}=1 \mid X_{i,t}\right)
  \;=\; \Lambda\!\left( \alpha_{a(i)} + X_{i,t}\,\beta_{s} \right),
  \qquad s \in \{1, 2, 4\}
\end{equation}

\noindent onde
\begin{itemize}
  \item $i$ é o hexágono H3 res 9, $t$ a semana, $a(i)$ a área FM do
        hexágono, $\alpha_{a(i)}$ o efeito fixo de área;
  \item $\Lambda(z) = e^{z}/(1+e^{z})$ é a função logística;
  \item $X_{i,t}$ é o vetor de regressores observados em $t$, e
        $\beta_s$ é o vetor de coeficientes estimado separadamente
        para cada horizonte $s$.
\end{itemize}

\noindent O estimador é um logit com regularização $L_2$
($C{=}1{,}0$, solver \texttt{liblinear}) e
\texttt{class\_weight='balanced'} para mitigar desbalanceamento de
classes (eventos raros). O split é \emph{temporal}: treino até
$T{-}4$, validação $T{-}3$ a $T{-}1$, teste em $T$ --- jamais
amostragem aleatória, dado o painel.

\paragraph{Grupos de regressores em $X_{i,t}$.} Extraídos da Fase 3 do
notebook:
\begin{enumerate}
  \item \textbf{Defasagens próprias do crime no hexágono} ---
        \texttt{y\_lag1}, \texttt{y\_lag2}, \texttt{y\_lag4},
        \texttt{y\_lag12} e a soma móvel \texttt{n\_crimes\_12w}.
        Capturam o mecanismo de \emph{near-repeat}.
  \item \textbf{Lag espacial} --- \texttt{W\_y\_lag1}, crime na
        vizinhança ring-1 do hexágono em $t{-}1$ \cite{JohnsonBowers2004}.
  \item \textbf{Sazonalidade harmônica} --- \texttt{week\_sin},
        \texttt{week\_cos} (ciclo anual) e \texttt{is\_holiday\_week}
        (semanas de feriado nacional/municipal).
  \item \textbf{Fatores urbanos} --- sete componentes
        (\texttt{n\_fat\_calcada}, \texttt{n\_fat\_ilum},
        \texttt{n\_fat\_lixo}, \texttt{n\_fat\_mobil},
        \texttt{n\_fat\_sitrua}, \texttt{n\_fat\_transito},
        \texttt{n\_fat\_vegetac}). Justifica-se manter cada componente
        em vez do total porque eles mapeiam diretamente nos órgãos
        responsáveis \cite{ObrienSampson2015, ChalfinEtAl2019}.
  \item \textbf{Cobertura por câmeras} --- \texttt{n\_cameras}
        (hex) e \texttt{n\_cameras\_ring1} (vizinhança).
  \item \textbf{Sinal qualitativo do Disque} ---
        \texttt{n\_dd\_lag1}, contagem de denúncias na semana anterior
        \cite{Wheeler2018}.
  \item \textbf{Contexto territorial ORCRIM} ---
        \texttt{orcrim\_CV}, \texttt{orcrim\_Mil\'icia},
        \texttt{orcrim\_TCP}, com a categoria
        \emph{``Sem domínio''} omitida como referência (ver caveat
        abaixo).
  \item \textbf{Efeitos fixos de área FM} ---
        \texttt{fe\_area\_*}, com uma área omitida como referência.
\end{enumerate}

\subsection{Como interpretar o modelo}
\label{ssec:pred_interpret}

A escolha de um logit linear nos regressores não é casual. Modelos mais
flexíveis (redes neurais, gradient boosting, random forest) podem
oferecer ganhos marginais de precisão, mas \emph{ao custo de uma caixa
preta}: o gestor recebe uma probabilidade prevista sem conseguir
explicar de onde ela veio. Aqui, a leitura é direta --- cada coeficiente
$\beta_k$ tem um significado próprio, em uma fórmula que cabe numa
linha. É essa transparência que sustenta o uso responsável da
ferramenta na decisão pública.

\paragraph{O que cada coeficiente significa.} No logit, o coeficiente
$\beta_k$ é o efeito de um aumento unitário no regressor $x_k$ sobre o
\emph{log-odds} de o crime ocorrer. Equivalentemente, $\exp(\beta_k)$ é
o \emph{odds ratio} --- o fator que multiplica a razão de chances de
$y{=}1$ contra $y{=}0$ quando $x_k$ aumenta em uma unidade, mantendo
todo o resto constante. Em linguagem natural:

\begin{itemize}
  \item $\exp(\beta_k) > 1$: a variável \emph{aumenta} a chance de crime.
  \item $\exp(\beta_k) < 1$: a variável \emph{reduz} a chance de crime.
  \item $\exp(\beta_k) \approx 1$: efeito praticamente neutro.
\end{itemize}

\paragraph{Exemplo isolado.} O coeficiente estimado para
\texttt{y\_lag1} (houve crime no mesmo hex na semana anterior) é
$\beta = +0{,}2346$. O odds ratio é $\exp(0{,}2346) \approx 1{,}26$.
\textbf{Leitura para o gestor:} um hex que teve crime na semana
passada enfrenta \emph{cerca de $26\%$ a mais de chance} de ter crime
na próxima, comparado a um hex idêntico em todo o resto que não teve.
Isto é a evidência quantitativa do mecanismo de \emph{near-repeat}
\cite{JohnsonBowers2004, MohlerEtAl2011} que a literatura já documenta.

\paragraph{Exemplo combinado: somando alavancas no log-odds.} Como o
preditor linear é uma soma, basta somar os $\beta$'s relevantes para
compor um cenário. Considere dois hexes hipotéticos, $A$ e $B$, na
mesma área FM (mesmo $\alpha$), na mesma semana (mesma sazonalidade).
A diferença está só nas variáveis abaixo:

\begin{table}[H]
  \centering
  \small
  \caption{Cenário ilustrativo: hex $A$ ``limpo'' versus hex $B$ com
  três alavancas adversas. Coeficientes $\beta$ extraídos da rodada
  atual de \texttt{coeficientes\_logit.csv} ($s{=}1$).}
  \label{tab:exemplo_interpret}
  \begin{tabular}{lrrrr}
    \toprule
    Variável & $\beta$ & $x_A$ & $x_B$ & Contribuição em $B$ \\
    \midrule
    \texttt{y\_lag1} (crime semana anterior)      & $+0{,}235$ & 0 & 1 & $+0{,}235$ \\
    \texttt{W\_y\_lag1} (crime vizinhança)        & $+0{,}366$ & 0 & 1 & $+0{,}366$ \\
    \texttt{n\_fat\_transito} (× 2 pontos)        & $+0{,}090$ & 0 & 2 & $+0{,}179$ \\
    \texttt{n\_fat\_calcada}                       & $+0{,}054$ & 0 & 1 & $+0{,}054$ \\
    \midrule
    \textbf{Diferença no log-odds ($B - A$)}      &            &   &   & $\mathbf{+0{,}834}$ \\
    \textbf{Odds ratio} $\exp(+0{,}834)$           &            &   &   & $\mathbf{\approx 2{,}30}$ \\
    \bottomrule
  \end{tabular}
\end{table}

\noindent\textbf{Tradução em risco absoluto.} Suponha que $A$ tem uma
probabilidade base $P(A) = 10\%$ de crime na próxima semana (próximo da
taxa observada no painel). Suas chances de crime são
$\text{odds}(A) = 0{,}10/0{,}90 \approx 0{,}111$. Para $B$, multiplicamos
pelo odds ratio: $\text{odds}(B) \approx 0{,}111 \times 2{,}30 \approx
0{,}255$. Convertendo de volta para probabilidade:
$P(B) = 0{,}255/(1+0{,}255) \approx 20\%$. \emph{Ou seja, $B$ tem cerca
do dobro da probabilidade absoluta de $A$, mesmo na mesma área e
semana.}

\paragraph{Como ler coeficientes negativos.} O coeficiente de
\texttt{orcrim\_CV} é $\beta = -1{,}90$, $\exp(-1{,}90) \approx
0{,}15$. \textbf{Leitura:} hexágonos em território de domínio CV têm,
em média no painel, cerca de $85\%$ \emph{menos} chance de roubo
registrado em comparação a hexes em ``Sem domínio'', controlando por
todo o resto. Isto \emph{não significa} que o CV reduz crime: é uma
correlação observada, possivelmente por subnotificação ou por efeito
inibidor sobre crime patrimonial em territórios sob controle de
facções. A leitura correta é \emph{``estar em CV está associado a
menos registros''}, nunca \emph{``CV causa menos crime''}.

\paragraph{Coeficientes desprezíveis também informam.} O coeficiente
de \texttt{n\_cameras} é $\beta = -0{,}003$, $\exp(-0{,}003) \approx
0{,}997$. \textbf{Leitura:} uma câmera adicional no hex está associada
a uma redução de cerca de $0{,}3\%$ nas chances --- praticamente
nada. Isto não invalida câmeras como política pública (elas geram
vídeo para investigação \emph{ex post}), mas mostra que, no painel
observado, sua presença não diferencia hexes de alto e baixo risco
\emph{ex ante}. O painel de variáveis acionáveis
(seção~\ref{ssec:acionaveis}) já sinaliza esta variável na categoria
\emph{Desprezível}.

\paragraph{Resumo prático para o gestor.}
\begin{enumerate}
  \item \emph{Sinal} do $\beta$: positivo aumenta risco; negativo reduz.
  \item \emph{Magnitude} de $\exp(\beta)$: distância de $1$ indica força.
        Acima de $1{,}5$ é forte; entre $1{,}05$ e $1{,}5$ é moderado;
        próximo de $1{,}00$ é desprezível.
  \item \emph{Efeitos se somam no log-odds e se multiplicam nas odds.}
        Combinações de fatores adversos podem dobrar ou triplicar a
        probabilidade rapidamente.
  \item \emph{Correlação, não causalidade.} Coeficientes orientam
        \emph{prioridade}, não garantem que intervir naquela variável
        reduzirá crime na mesma proporção.
\end{enumerate}

\subsection{Métricas de validação}
\label{ssec:pred_metrics}

O modelo é avaliado em quatro métricas, computadas sobre o conjunto de
teste para cada horizonte de previsão:

\begin{itemize}
  \item \textbf{AUC-ROC} --- discriminação ordinal entre hexes que
        terão e que não terão crime na semana $t{+}s$.
  \item \textbf{AUC-PR} --- precisão-recall, mais sensível ao
        desbalanceamento de classes (eventos raros).
  \item \textbf{PAI top10 (Prediction Accuracy Index)} ---
        razão entre a fração de eventos capturada nos 10\% de hexes de
        maior risco previsto e a fração de área coberta por esses 10\%- \textbf{É a principal métrica para policy, na opinião dos autores, porque reflete a eficiência da alocação de recursos: quanto mais eventos capturados em menos área, melhor.}
        Métrica padrão na literatura de \emph{predictive policing}
        \cite{ChaineyEtAl2008}.
  \item \textbf{Moran's I dos resíduos} --- autocorrelação espacial
        residual; valores próximos de zero indicam que o modelo
        capturou a estrutura espacial do crime.
\end{itemize}

\noindent Resultados atuais (extraídos de
\texttt{metricas\_validacao.csv}):

\begin{table}[H]
  \centering
  \caption{Métricas de validação por horizonte de previsão. Conjunto de
  teste com 156 observações em todos os horizontes; conjunto de treino
  com 38{.}220 observações.}
  \label{tab:metrics}
  \begin{tabular}{lrrrr}
    \toprule
    Horizonte & AUC-ROC & AUC-PR & PAI top10 & Moran's I ($p$) \\
    \midrule
    $s{=}1$ (1 sem.)  & 0{,}7275 & 0{,}3713 & 3{,}5294 & 0{,}1350 ($p{=}0{,}014$) \\
    $s{=}2$ (2 sem.)  & 0{,}7529 & 0{,}3871 & 3{,}5294 & 0{,}1167 ($p{=}0{,}017$) \\
    $s{=}4$ (4 sem.)  & 0{,}7088 & 0{,}3759 & 3{,}5294 & 0{,}1434 ($p{=}0{,}005$) \\
    \bottomrule
  \end{tabular}
\end{table}

\begin{figure}[H]
  \centering
  \figfile{17_metricas_tab.png}{0.7\linewidth}
  \caption{Aba ``Métricas'' do Mapa Preditivo.}
  \label{fig:metricas}
\end{figure}

\subsection{Drivers por área}
\label{ssec:drivers}

A aba \emph{Drivers} mostra, para cada área FM, os cinco hexágonos com
maior probabilidade prevista e os três fatores que mais contribuem para
o risco previsto em cada um (top contribuintes $|\beta \cdot x|$).

\paragraph{Cuidado de leitura.} O \emph{rank} ordena os hexágonos
\emph{dentro de cada área FM}, não no município inteiro. Isto garante
que todas as áreas apareçam representadas; em contrapartida, P(crime)
entre áreas diferentes não é diretamente comparável --- um \emph{rank
1} em uma área pode ter risco absoluto menor que um \emph{rank 5} em
outra. A própria tela exibe esta nota acima da tabela.

\begin{figure}[H]
  \centering
  \figfile{16_drivers_tab.png}{0.95\linewidth}
  \caption{Aba ``Drivers'' --- top hexes por área e seus contribuintes.}
  \label{fig:drivers}
\end{figure}

\subsection{Painel de variáveis acionáveis}
\label{ssec:acionaveis}

Esta é a \textbf{contribuição central} do painel preditivo para o
gestor. Os regressores do modelo se dividem em duas naturezas:
\begin{itemize}
  \item \textbf{Variáveis estruturais} --- defasagens do próprio
        crime, lag espacial, efeitos fixos de área, sazonalidade,
        domínio de facção. São úteis para \emph{prever crime}, mas
        \textbf{não orientam intervenção} (o gestor não pode reduzir
        a sazonalidade nem o passado).
  \item \textbf{Alavancas de política pública} --- fatores urbanos
        desagregados, câmeras, denúncias do Disque. \emph{O município
        pode agir sobre estas.}
\end{itemize}

\noindent O painel filtra os coeficientes para mostrar
\emph{somente} as alavancas, com:
direção do efeito ($\uparrow$ aumenta risco / $\downarrow$ reduz /
$\approx$ neutro),
\emph{magnitude qualitativa} pelo valor absoluto de $\beta$
(\emph{Notável} $|\beta| > 0{,}05$; \emph{Leve}
$0{,}01 \le |\beta| \le 0{,}05$; \emph{Desprezível}
$|\beta| < 0{,}01$),
órgãos responsáveis e uma nota curta.

\paragraph{Caveat causal.} Os coeficientes vêm de um modelo
\emph{preditivo}, não de um experimento causal: a correlação observada
justifica priorização, não causalidade. Para câmeras, há ainda
\emph{viés de seleção} (instalação endógena onde já há crime, o que
gera correlação positiva com risco mesmo quando o efeito causal pode
ser nulo ou negativo) \cite{LumIsaac2016, EnsignEtAl2018}.

\begin{figure}[H]
  \centering
  \figfile{19_painel_alavancas.png}{0.95\linewidth}
  \caption{Painel de variáveis acionáveis. As alavancas estão
  ordenadas por $|\beta_{s=1}|$ decrescente.}
  \label{fig:alavancas}
\end{figure}

\subsection{Mapa interativo e popups por hexágono}
\label{ssec:mapa_interativo}

O mapa interativo serve as predições T+1 por padrão, com toggles para
T+2 e T+4 e overlay dos polígonos FM. Cada hexágono, ao ser clicado,
abre um popup com:
área(s) FM, horizonte, P(crime), decil de risco e os três fatores que
mais explicam aquele risco.

\begin{figure}[H]
  \centering
  \figfile{14_mapa_preditivo_geral.png}{0.95\linewidth}
  \caption{Mapa interativo --- horizonte T+1.}
  \label{fig:mapa_t1}
\end{figure}

\begin{figure}[H]
  \centering
  \figfile{15_mapa_preditivo_area_PV.png}{0.95\linewidth}
  \caption{Mapa preditivo filtrado por área (Presidente Vargas ---
  Campo de Santana --- Central do Brasil --- Cinelândia), mostrando o
  recorte de hexágonos dentro do polígono da Força Municipal.}
  \label{fig:mapa_polig}
\end{figure}

\begin{figure}[H]
  \centering
  \figfile{20_popup_hex.png}{0.7\linewidth}
  \caption{Popup detalhado de um hexágono, com os três principais
  drivers.}
  \label{fig:popup}
\end{figure}

\subsection{Coeficientes em português}
\label{ssec:coef}

A aba \emph{Coeficientes} mostra todos os regressores estimados com
seus \emph{odds ratios} por horizonte ($\exp(\beta)$). Para reduzir
fricção de leitura, os nomes técnicos das variáveis são traduzidos
para rótulos em português natural (p.\ ex.\ \texttt{y\_lag1}
$\to$ ``Crime no mesmo trecho (1 semana antes)'';
\texttt{n\_fat\_ilum} $\to$ ``Fator urbano: iluminação''). O nome
técnico continua sendo a chave canônica nos CSVs.

\begin{figure}[H]
  \centering
  \figfile{18_coeficientes_tab.png}{0.95\linewidth}
  \caption{Aba ``Coeficientes'' com rótulos em português.}
  \label{fig:coef}
\end{figure}

% =============================================================================
\section{Guardrails de IA responsável}
\label{sec:guardrails}

Este é um projeto de segurança pública. As regras abaixo são
\emph{inegociáveis} e aparecem em todas as superfícies da aplicação.

\begin{enumerate}
  \item \textbf{Decisão final sempre humana.} A plataforma produz
        \emph{rascunhos} auditáveis (relatório, sugestões, escores). O
        Relatório exportado em DOCX vem marcado com banner
        \emph{``RASCUNHO --- validação humana''}.
  \item \textbf{Foco no ambiente, não no indivíduo.} A plataforma
        analisa fatores urbanos e alocação de patrulha, não vigia
        pessoas \cite{CohenFelson1979, ObrienSampson2015}.
  \item \textbf{Texto livre é indício, não fato.} Disque Denúncia e
        RELINTs sempre vêm com citação literal, fonte rastreável e
        nível de confiança. A IA sinaliza incerteza em vez de afirmar.
  \item \textbf{Nunca inventar dado.} Camada ausente é declarada como
        ausente. Escore com dado insuficiente é marcado como
        \emph{baixa confiança}.
  \item \textbf{Cuidado com o viés de realimentação.} Priorizar sempre
        as mesmas áreas gera mais patrulha, mais registro e escore
        artificialmente mais alto na próxima iteração. A literatura
        documenta este risco em modelos preditivos de policiamento
        \cite{LumIsaac2016, EnsignEtAl2018}; a plataforma deve ser
        usada com revisão humana periódica.
  \item \textbf{LGPD.} Dados são tratados como agregados
        territoriais. Dados pessoais (Disque, CPSR) são minimizados.
        O copiloto não tem acesso a SQL livre --- apenas a uma
        \emph{whitelist} de consultas.
  \item \textbf{Transparência da fórmula.} O escore da S10 é uma soma
        ponderada com pesos extraídos em constantes nomeadas; a fórmula
        é exibida na própria tela. Não há caixa-preta.
\end{enumerate}

% =============================================================================
\begin{thebibliography}{99}

\bibitem{ChaineyEtAl2008}
Chainey, S., Tompson, L., \& Uhlig, S. (2008).
The Utility of Hotspot Mapping for Predicting Spatial Patterns of Crime.
\emph{Security Journal}, 21(1--2), 4--28.

\bibitem{MohlerEtAl2011}
Mohler, G.~O., Short, M.~B., Brantingham, P.~J., Schoenberg, F.~P.,
\& Tita, G.~E. (2011).
Self-Exciting Point Process Modeling of Crime.
\emph{Journal of the American Statistical Association}, 106(493), 100--108.

\bibitem{MohlerEtAl2015}
Mohler, G.~O., Short, M.~B., Malinowski, S., Johnson, M., Tita, G.~E.,
Bertozzi, A.~L., \& Brantingham, P.~J. (2015).
Randomized Controlled Field Trials of Predictive Policing.
\emph{Journal of the American Statistical Association}, 110(512), 1399--1411.

\bibitem{JohnsonBowers2004}
Johnson, S.~D., \& Bowers, K.~J. (2004).
The Burglary as Clue to the Future: The Beginnings of Prospective
Hot-Spotting. \emph{European Journal of Criminology}, 1(2), 237--255.

\bibitem{CohenFelson1979}
Cohen, L.~E., \& Felson, M. (1979).
Social Change and Crime Rate Trends: A Routine Activity Approach.
\emph{American Sociological Review}, 44(4), 588--608.

\bibitem{Wheeler2018}
Wheeler, A.~P. (2018).
The Effect of 311 Calls for Service on Crime in Washington, DC, at
Microplaces. \emph{Crime \& Delinquency}, 64(14), 1882--1903.

\bibitem{ObrienSampson2015}
O'Brien, D.~T., \& Sampson, R.~J. (2015).
Public and Private Spheres of Neighborhood Disorder: Assessing
Pathways to Violence Using Large-Scale Digital Records.
\emph{Journal of Research in Crime and Delinquency}, 52(4), 486--510.

\bibitem{CaplanKennedyMiller2011}
Caplan, J.~M., Kennedy, L.~W., \& Miller, J. (2011).
Risk Terrain Modeling: Brokering Criminological Theory and GIS Methods
for Crime Forecasting. \emph{Justice Quarterly}, 28(2), 360--381.

\bibitem{ChalfinEtAl2019}
Chalfin, A., Hansen, B., Lerner, J., \& Parker, L. (2019).
Reducing Crime through Environmental Design: Evidence from a
Randomized Experiment of Street Lighting in New York City.
\emph{NBER Working Paper} No.\ 25798.

\bibitem{LumIsaac2016}
Lum, K., \& Isaac, W. (2016).
To Predict and Serve? \emph{Significance}, 13(5), 14--19.

\bibitem{EnsignEtAl2018}
Ensign, D., Friedler, S.~A., Neville, S., Scheidegger, C.,
\& Venkatasubramanian, S. (2018).
Runaway Feedback Loops in Predictive Policing.
\emph{Proceedings of Machine Learning Research (FAT*)}, 81, 160--171.

\end{thebibliography}

\end{document}
